\documentclass[11pt,a4paper]{article}

% ---------- Packages ----------
\usepackage[utf8]{inputenc}
\usepackage[T1]{fontenc}
\usepackage[margin=2.5cm]{geometry}
\usepackage{amsmath}
\usepackage{booktabs}
\usepackage{array}
\usepackage{xcolor}
\usepackage{listings}
\usepackage{fancyvrb}
\usepackage{hyperref}
\usepackage{enumitem}
\setlist{topsep=2pt, itemsep=1pt, parsep=0pt}

% ---------- Hyperlinks ----------
\hypersetup{
  colorlinks=true,
  linkcolor=blue!60!black,
  urlcolor=blue!60!black,
  citecolor=blue!60!black,
  pdftitle={Mini RAG Project Documentation},
  pdfauthor={Project Author}
}

% ---------- Listings style ----------
\definecolor{codebg}{RGB}{248,248,248}
\definecolor{codekeword}{RGB}{62,88,151}
\definecolor{codecomment}{RGB}{96,140,98}
\definecolor{codestring}{RGB}{163,23,23}

\lstdefinestyle{pythonstyle}{
  language=Python,
  basicstyle=\ttfamily\small,
  keywordstyle=\color{codekeword}\bfseries,
  commentstyle=\color{codecomment}\itshape,
  stringstyle=\color{codestring},
  backgroundcolor=\color{codebg},
  frame=single,
  rulecolor=\color{gray!40},
  breaklines=true,
  showstringspaces=false,
  tabsize=4,
  numbers=left,
  numberstyle=\tiny\color{gray},
  numbersep=6pt,
  xleftmargin=14pt,
  framexleftmargin=14pt,
  captionpos=b
}

\lstdefinestyle{bashstyle}{
  language=bash,
  basicstyle=\ttfamily\small,
  backgroundcolor=\color{codebg},
  frame=single,
  rulecolor=\color{gray!40},
  breaklines=true,
  showstringspaces=false,
  tabsize=2,
  captionpos=b
}

\lstset{style=pythonstyle}

% ---------- Inline code ----------
\newcommand{\code}[1]{\texttt{#1}}

% ---------- Title ----------
\title{\textbf{Mini RAG: A Minimal From-Scratch\\ Retrieval-Augmented Generation Pipeline}\\[4pt]
\large Project Documentation}
\author{Project Author}
\date{\today}

\begin{document}

\maketitle

\begin{center}
\textit{Keywords: Retrieval-Augmented Generation, Sentence Embeddings,
Cosine Similarity, FastAPI, Google Gemini, Amazon Bedrock}
\end{center}

\begin{abstract}
\noindent
This document provides complete technical documentation for a minimal,
from-scratch Retrieval-Augmented Generation (RAG) pipeline. The project
ingests documents in PDF, plain-text, and Markdown formats as well as web
articles, splits them into overlapping word-based chunks, embeds each chunk
locally using a sentence-transformer model, and serves similarity-based
retrieval over a plain \texttt{numpy} vector store. A generation step then
produces answers that are grounded exclusively in the retrieved context,
using either Google Gemini or Amazon Bedrock (Claude) as an interchangeable
runtime backend. The system exposes both a browser-based user interface and a
command-line interface. Every stage---chunking, embedding, storage,
retrieval, and generation---is implemented in readable, dependency-light
Python with no orchestration framework (such as LangChain) and no dedicated
vector-database server, making the pipeline fully transparent and easy to
study, tune, or replace piece by piece. This report covers the system
architecture, module-by-module reference, configuration, API specification,
persistence model, containerization, deployment paths, security
considerations, and known limitations.
\end{abstract}

\newpage
\tableofcontents
\newpage

% =================================================================
\section{Introduction}

\subsection{What Is Retrieval-Augmented Generation?}
Retrieval-Augmented Generation (RAG) is a pattern in which a language model
answers a question using a restricted set of passages retrieved from a
user-supplied corpus, rather than relying solely on its parametric
knowledge. A RAG pipeline consists of three principal stages:
\begin{enumerate}
  \item \textbf{Ingestion and chunking}: documents are loaded and split into
        smaller, self-contained text fragments.
  \item \textbf{Embedding and retrieval}: each fragment is converted into a
        dense vector; at query time the question is embedded with the same
        model and the most similar fragments are returned.
  \item \textbf{Grounded generation}: the retrieved fragments are inserted
        into a prompt and an LLM produces an answer using only that context.
\end{enumerate}

\subsection{Project Goals}
This project implements the full RAG loop from scratch with the explicit
intention of \emph{seeing every step}. The design decisions follow from this
goal:
\begin{itemize}
  \item \textbf{No orchestration framework} such as LangChain; the pipeline
        is plain Python.
  \item \textbf{No vector database server}; a \texttt{numpy} array performs
        the similarity search and persists to a pickle file
        (\code{index.pkl}).
  \item \textbf{Local, free embeddings}; the
        \texttt{all-MiniLM-L6-v2} sentence-transformer model requires no API
        key.
  \item \textbf{Swappable LLM backends}; a single environment variable
        switches generation between Google Gemini and Amazon Bedrock
        (Claude) at runtime.
  \item \textbf{Two interaction surfaces}: a browser UI served by FastAPI
        and a set of small command-line scripts.
\end{itemize}

% =================================================================
\section{System Architecture}

\subsection{High-Level Data Flow}
Both the browser upload panel and the command-line entry points funnel into
the same chunk $\rightarrow$ embed $\rightarrow$ store pipeline; ingestion is
the only layer that is aware of file formats or URLs. Figure~\ref{fig:pipeline}
illustrates the flow.

\begin{figure}[h]
\centering
\begin{Verbatim}[frame=single, fontsize=\small, numbers=none]
 Browser upload / URL            your_docs/*.md  (CLI path)
        |                              |
        v                              v
   POST /ingest/file |         build_chunks(folder)
   POST /ingest/url  +----------------> rag/ingest.py
                                             |  splits each source into
                                             v  overlapping word chunks
                                       rag/store.py
                                             |  embeds each chunk with
                                             |  all-MiniLM-L6-v2, keeps
                                             |  vectors in numpy
                                             v
                          POST /query OR query.py
                                             |  embeds the question, finds
                                             |  the top-3 most similar
                                             |  chunks (cosine similarity),
                                             |  builds a grounded prompt
                                             v
                                      rag/generate.py
                                             |  generates an answer via
                                             |  Gemini or Bedrock
                                             v
                                         answer text
\end{Verbatim}
\caption{Pipeline data flow.}
\label{fig:pipeline}
\end{figure}

\subsection{Repository Layout}
Table~\ref{tab:repo} lists every file in the repository and its role.

\begin{table}[h]
\centering
\small
\begin{tabular}{>{\ttfamily}p{4.4cm} p{10.4cm}}
\toprule
\multicolumn{1}{>{\rmfamily}p{4.4cm}}{\textbf{File}} &
\multicolumn{1}{c}{\textbf{Role}} \\
\midrule
api.py            & FastAPI server: REST endpoints plus browser-UI hosting \\
build\_index.py   & CLI: chunk, embed, and persist an entire folder of documents \\
query.py          & CLI: ask a question from the terminal \\
fetch\_url.py     & CLI: download a URL's article text into \texttt{sample\_docs/} \\
requirements.txt  & Python dependencies with pinned minimum versions \\
Dockerfile        & Container image definition \\
.dockerignore     & Files excluded from the Docker image \\
.env              & Runtime configuration (gitignored; never committed) \\
.gitignore        & Files ignored by git \\
index.pkl         & Persisted vector index (pickle format; gitignored) \\
README.md         & Project documentation (this document's source) \\
DEPLOY\_AWS.md    & EC2 + venv + Gemini deployment guide \\
DEPLOY\_DOCKER\_BEDROCK.md & EC2 + Docker + Bedrock deployment guide \\
rag/\_\_init\_\_.py & Empty package marker \\
rag/ingest.py     & Document loading and chunking \\
rag/store.py      & Embedding, similarity search, and persistence \\
rag/generate.py   & LLM generation (Gemini / Bedrock backends) \\
static/index.html & Single-page frontend (no build step, no framework) \\
sample\_docs/     & Sample documents for the CLI path \\
\bottomrule
\end{tabular}
\caption{Repository layout.}
\label{tab:repo}
\end{table}

% =================================================================
\section{Technology Stack}
Table~\ref{tab:stack} summarizes the tools used and their role in the
project.

\begin{table}[h]
\centering
\small
\begin{tabular}{>{\raggedright}p{3.4cm} >{\raggedright}p{4.6cm} p{7.0cm}}
\toprule
\textbf{Tool} & \textbf{Category} & \textbf{Function in this project} \\
\midrule
Python $\geq$ 3.11 & Language & Implementation language; the shipped Docker image uses 3.11-slim \\
FastAPI + uvicorn & Web framework & REST API endpoints; uvicorn is the ASGI server running the app \\
sentence-transformers & Embedding model & Converts chunks and questions into 384-dimensional dense vectors with \texttt{all-MiniLM-L6-v2} \\
numpy & Math / storage & Stores all embeddings in one matrix; implements the cosine-similarity search as a dot product; persists via pickle \\
pypdf & PDF parsing & Extracts raw text, page by page, from text-layer PDFs \\
trafilatura & Web scraping & Fetches a URL and extracts article text, title, and metadata; strips navigation and ads \\
google-genai & LLM client & Calls Google Gemini (\texttt{generate\_content}) when $GENERATION\_BACKEND=gemini$ \\
boto3 & AWS client & Calls the Amazon Bedrock Converse API (Claude) when $GENERATION\_BACKEND=bedrock$; auto-uses IAM role credentials \\
python-dotenv & Configuration & Loads \code{.env} into environment variables at startup \\
Docker & Containerization & Reproducible build/deploy of the app (Docker + Bedrock path) \\
AWS EC2, IAM, systemd & Deployment & Hosts the app on Ubuntu; systemd keeps uvicorn alive; IAM provides Bedrock access without keys \\
Vanilla HTML/CSS/JS & Frontend & Single-file UI: upload/drag-drop, URL paste, source list, question box, answer rendering \\
\bottomrule
\end{tabular}
\caption{Technology stack.}
\label{tab:stack}
\end{table}

% =================================================================
\section{The RAG Pipeline in Detail}

\subsection{Ingestion (\code{rag/ingest.py})}
The ingestion layer is format-aware but the rest of the pipeline is not. Its
public symbols are:

\begin{itemize}
  \item \code{Chunk} --- a dataclass with fields \code{text} (the fragment),
        \code{source} (filename the chunk came from), and \code{chunk\_id}
        (zero-based position within the file).
  \item \code{extract\_pdf\_text(path)} --- opens a PDF with
        \texttt{pypdf.PdfReader} and joins every page's extraction with a
        blank line. Pages without a text layer yield \code{""} rather than an
        exception.
  \item \code{slugify(text, max\_len=60)} --- builds a safe filename slug
        from a title: lowercases, collapses runs of non-alphanumeric
        characters into hyphens, truncates to \code{max\_len}, and falls
        back to \code{"page"} if empty.
  \item \code{extract\_url\_text(url)} --- downloads a page with
        \texttt{trafilatura.fetch\_url} and extracts the readable article
        (comments and tables excluded). Raises \code{ValueError} if the page
        cannot be downloaded or contains no readable content. Returns a
        \code{(title, text)} tuple.
  \item \code{load\_documents(folder)} --- returns a list of
        \code{(filename, text)} pairs for every \code{.md}, \code{.txt}, and
        \code{.pdf} file in a folder (sorted by filename); other files are
        skipped silently.
  \item \code{chunk\_text(text, chunk\_size=500, overlap=50)} --- word-based
        chunking: splits on whitespace and steps forward \code{chunk\_size}
        words while re-including the previous \code{overlap} words so that
        sentences straddling a boundary keep context.
  \item \code{chunk\_single\_document(name, text, ...)} --- wraps the output
        of \code{chunk\_text} into \code{Chunk} objects; used both by the
        folder path and by the API's incremental ingest endpoints.
  \item \code{build\_chunks(folder, chunk\_size=500, overlap=50)} --- loads
        every document in a folder and flattens all chunks into one list.
\end{itemize}

A module-level sanity check is available via \code{python -m rag.ingest},
which prints the chunk count and a preview of the first two chunks.

\begin{table}[h]
\centering
\small
\begin{tabular}{lccl}
\toprule
\textbf{Entry point} & \textbf{chunk\_size} & \textbf{overlap} & \textbf{Chunk length} \\
\midrule
CLI (\code{build\_index.py})            & 300 & 50 & $\approx$ 300 words \\
Browser/API (\code{api.py} / defaults)  & 500 & 50 & $\approx$ 500 words \\
\bottomrule
\end{tabular}
\caption{Chunk-size discrepancy between the two entry points.}
\label{tab:chunk}
\end{table}

As Table~\ref{tab:chunk} highlights, the CLI path and the API path use
different default chunk sizes; both are internally consistent but produce
different chunk counts.

\subsection{Embedding and Storage (\code{rag/store.py})}
\begin{itemize}
  \item \code{EMBEDDING\_MODEL = "all-MiniLM-L6-v2"} --- a small, fast,
        free sentence-transformer producing \textbf{384-dimensional},
        \textbf{normalized} vectors.
  \item \code{VectorStore} holds a list of \code{Chunk}s and a
        \texttt{numpy} embedding matrix (\code{embeddings}).
  \item \code{add(chunks)} --- encodes new texts with
        \code{normalize\_embeddings=True} and appends them, growing the
        matrix with \code{np.vstack}.
  \item \code{search(query, top\_k=3)} --- encodes the query and computes
        \textbf{cosine similarity} as a matrix--vector dot product (valid
        because all vectors are normalized). Returns the \code{top\_k}
        highest-scoring results in descending order.
  \item \code{sources()} --- returns the deduplicated list of source
        filenames in insertion (first-seen) order.
  \item \code{remove\_source(source)} --- deletes every chunk and embedding
        row belonging to a source.
  \item \code{save(path)} / \code{load(path)} --- pickles or unpickles the
        dictionary \code{\{"chunks": [Chunk, ...], "embeddings": ndarray\}}.
\end{itemize}

\subsection{Generation (\code{rag/generate.py})}
\code{generate\_answer(prompt)} reads the environment variable
\code{GENERATION\_BACKEND} (default \code{"gemini"}) and dispatches to one of
two providers. Both \code{api.py} and \code{query.py} call exactly this one
function, so the LLM provider can be swapped without touching the rest of the
pipeline.

\begin{itemize}
  \item \textbf{Gemini backend}: lazily imports \code{google.genai},
        requires \code{GEMINI\_API\_KEY}, and calls
        \code{client.models.generate\_content} with the model named by
        \code{GEMINI\_MODEL} (default \texttt{gemini-2.5-flash}).
  \item \textbf{Bedrock backend}: lazily imports \code{boto3} and calls the
        provider-agnostic Converse API
        (\code{client.converse}) with an inference cap of 500 tokens
        (\code{maxTokens}). Credentials are resolved automatically by
        \texttt{boto3} from environment variables, \code{\textasciitilde{}/.aws/credentials},
        or an EC2 IAM instance role, so no secret material needs to live in
        \code{.env} when deployed on EC2.
\end{itemize}

% =================================================================
\section{Application Layer}

\subsection{REST Server (\code{api.py})}
The FastAPI application is created at module scope; the \code{VectorStore}
and the embedding model are loaded \textbf{once at startup} rather than per
request, avoiding the multi-second model-load cost on every query. A missing
\code{index.pkl} is tolerated so the first document can be added from the
browser. Table~\ref{tab:api} specifies the endpoints; the full request and
response schemas, including error bodies, are given in
Section~\ref{sec:api-ref}.

\begin{table}[h]
\centering
\small
\begin{tabular}{>{\ttfamily}p{2.6cm} p{4.8cm} p{7.4cm}}
\toprule
\textbf{Method/Path} & \textbf{Request} & \textbf{Behavior and response} \\
\midrule
GET /                & --- & Serves the single-page frontend \\
GET /sources         & --- & List of distinct indexed filenames plus total chunk count \\
POST /ingest/file    & multipart file & Saves, chunks, embeds, persists; returns filename, chunks added, total \\
POST /ingest/url     & JSON \texttt{\{url\}} & Fetches article, writes file, chunks, embeds, persists \\
POST /query          & JSON \texttt{\{question\}} & Retrieves top-3 chunks, generates a grounded answer \\
DELETE /sources/\{filename\} & --- & Removes one source and all its chunks; 404 if not indexed \\
\bottomrule
\end{tabular}
\caption{REST API endpoints.}
\label{tab:api}
\end{table}

\subsection{Frontend (\code{static/index.html})}
The interface is a single HTML file with inline CSS and vanilla JavaScript:
\begin{itemize}
  \item A \emph{source panel} with two tabs, \emph{Upload file} (click or
        drag-and-drop, accepts \code{.pdf}, \code{.txt}, \code{.md}) and
        \emph{Paste URL}; an ingest-status line; and a list of indexed
        sources, each with a Remove button.
  \item A \emph{question bar} and answer area with a pulsing
        \emph{Searching \& thinking} indicator while a request is in flight.
  \item All dynamically rendered text (source names, answers) is passed
        through an \code{escapeHtml} helper (a \code{div.textContent} trick)
        before being inserted, preventing HTML injection.
  \item A \code{prefers-reduced-motion} media query disables the loading-dot
        animation for users who request reduced motion.
\end{itemize}

\subsection{Command-Line Tools}
\begin{itemize}
  \item \code{build\_index.py} --- accepts an optional folder argument
        (default \code{sample\_docs/}), chunks with
        \code{chunk\_size=300, overlap=50}, embeds, and persists to
        \code{index.pkl}. It \emph{rebuilds from scratch} rather than
        merging with an existing index.
  \item \code{query.py} --- loads the store, retrieves the top-3 chunks,
        prints each as \code{[score:.3f] source (chunk id)}, and prints the
        generated answer.
  \item \code{fetch\_url.py} --- downloads a URL and writes
        \code{\# title + Source: url + article text} into
        \code{sample\_docs/<slug>.txt} for \code{build\_index.py} to pick up.
\end{itemize}

% =================================================================
\section{Installation and Usage}

\subsection{Local Setup}
\begin{lstlisting}[language=bash, style=bashstyle]
python -m venv .venv
# Windows:     .venv\Scripts\activate
# macOS/Linux: source .venv/bin/activate
pip install -r requirements.txt
\end{lstlisting}

The first embedding run downloads \texttt{all-MiniLM-L6-v2} from Hugging
Face (approximately 90\,MB). The generation backend and its key/model must be
configured in \code{.env} before questions can be answered.

\subsection{Running the Web Application}
\begin{lstlisting}[language=bash, style=bashstyle]
uvicorn api:app --reload
\end{lstlisting}
Open \textbf{http://localhost:8000}. Add a source (file upload or URL), then
ask a question grounded only in the added sources.

\subsection{Command-Line Workflow}
\begin{lstlisting}[language=bash, style=bashstyle]
python build_index.py path/to/your/folder   # chunk + embed a folder
python query.py "your question"              # ask from the terminal
python fetch_url.py https://example.com/post # save an article
\end{lstlisting}

% =================================================================
\section{Configuration}
\code{.env} is read by \code{load\_dotenv()} in both \code{api.py} and
\code{query.py} and is gitignored. Every variable has a code-level default.
Table~\ref{tab:env} summarizes all settings.

\begin{table}[h]
\centering
\small
\begin{tabular}{>{\ttfamily}p{3.5cm} >{\ttfamily}p{3.6cm} p{7.7cm}}
\toprule
\textbf{Variable} & \textbf{Code default} & \textbf{Meaning} \\
\midrule
INDEX\_PATH             & index.pkl & Pickle file for the vector index \\
DOCS\_FOLDER            & sample\_docs & Where uploads and URL extractions are stored \\
GENERATION\_BACKEND     & gemini    & \code{gemini} or \code{bedrock} \\
GEMINI\_API\_KEY        & (required) & Google AI Studio key for the Gemini backend \\
GEMINI\_MODEL           & gemini-2.5-flash & Gemini model name \\
AWS\_REGION             & us-east-1 & AWS region for Bedrock \\
BEDROCK\_MODEL\_ID      & anthropic.claude-3-5-haiku-20241022-v1:0 & Bedrock model identifier \\
\bottomrule
\end{tabular}
\caption{Environment-variable configuration.}
\label{tab:env}
\end{table}

\begin{table}[h]
\centering
\small
\begin{tabular}{>{\ttfamily}p{3.5cm} p{11.3cm}}
\toprule
\textbf{Variable} & \textbf{Value shipped in local .env} \\
\midrule
GENERATION\_BACKEND & gemini \\
GEMINI\_MODEL       & gemini-2.5-flash \\
AWS\_REGION         & ap-south-2 \\
BEDROCK\_MODEL\_ID  & anthropic.claude-haiku-4-5-20251001-v1:0 \\
\bottomrule
\end{tabular}
\caption{Variables actually set in the provided \code{.env}.}
\label{tab:env-shipped}
\end{table}

\textbf{Heads-up}: several scripts and deployment guides reference a
\code{.env.example} file that is \emph{not} present in the repository;
a copy of \code{.env} or a manually created \code{.env.example} is needed
before following the \code{cp .env.example .env} steps in the deployment
documents.

% =================================================================
\section{REST API Reference}\label{sec:api-ref}

\subsection{GET /sources}
Response body:
\begin{lstlisting}[language=Python, style=pythonstyle]
{"sources": [<filename>, ...], "total_chunks": <int>}
\end{lstlisting}
Returns the deduplicated, insertion-ordered list of indexed filenames and the
total number of chunks. Empty store yields empty list and zero.

\subsection{POST /ingest/file}
Multipart form field \code{file}; accepts \code{.pdf}, \code{.txt},
\code{.md}.
\begin{itemize}
  \item 400 --- unsupported extension: \emph{Unsupported file type `ext'. Use
        .pdf, .txt, or .md.}
  \item 400 --- no extractable text: \emph{No extractable text found in that
        file. If it's a scanned PDF, it needs OCR first.}
\end{itemize}
On success the file is written into \code{DOCS\_FOLDER}, chunked, embedded,
and persisted. Response:
\begin{lstlisting}[language=Python, style=pythonstyle]
{"filename": <name>, "chunks_added": <int>, "total_chunks": <int>}
\end{lstlisting}

\subsection{POST /ingest/url}
JSON body \code{\{"url": "https://..."\}}. Raises 400 with the
\code{ValueError} message from \code{extract\_url\_text} on failure. Writes
\code{\# <title> / Source: <url> / <text>} to
\code{DOCS\_FOLDER/<slug>.txt}, then chunk/embed/persist; response shape as
above.

\subsection{POST /query}
JSON body \code{\{"question": "..."\}}.
\begin{itemize}
  \item Empty store $\Rightarrow$ 200 with answer \emph{Nothing's been
        indexed yet --- add a file or URL above first.}
  \item Retrieves top-3 chunks, builds the grounded prompt, calls
        \code{generate\_answer}.
  \item Generation failure $\Rightarrow$ 200 with
        \code{"ERROR generating answer: \{e\}"} (errors are returned as
        answers, not as HTTP errors).
\end{itemize}
Response: \code{\{"answer": "<generated text>"\}}.

\subsection{DELETE /sources/\{filename\}}
Removes one source and its chunks/embeddings, persists, and returns
\code{\{"removed": <filename>, "total_chunks": <int>\}}. 404 with
\emph{'\texttt{filename}' isn't currently indexed.} if absent.

% =================================================================
\section{Data Persistence}
\begin{itemize}
  \item \textbf{Index file} (\code{index.pkl}): a pickle of
        \code{\{"chunks": [Chunk, ...], "embeddings": ndarray\}}, saved after
        every API ingest/delete and by \code{build\_index.py}; loaded once at
        API startup. Gitignored and excluded from Docker images.
  \item \textbf{Document folder} (\code{DOCS\_FOLDER}, default
        \code{sample\_docs/}): uploaded files are stored verbatim; URL
        ingests write \code{\# title / Source: url / text} files there. This
        folder is \emph{not} re-ingested automatically at startup --- the
        index is self-contained.
  \item \textbf{Memory}: chunks and embeddings live in RAM (numpy arrays)
        for the process lifetime.
  \item \textbf{Concurrency}: there is no locking around \code{store.add} /
        \code{store.save}; concurrent ingests could in principle race on the
        pickle write. Acceptable for single-user use.
\end{itemize}

% =================================================================
\section{Containerization (Docker)}
The \code{Dockerfile} is based on \code{python:3.11-slim}:
\begin{enumerate}
  \item \code{WORKDIR /app}.
  \item Dependencies are installed first (separate cacheable layer) with
        \code{pip install --no-cache-dir -r requirements.txt}, so the heavy
        torch/transformers layer is rebuilt only when requirements change.
  \item The application code is copied in.
  \item \code{RUN mkdir -p sample\_docs} ensures the upload folder exists.
  \item \code{EXPOSE 8000}; \code{CMD ["uvicorn", "api:app", "--host",
        "0.0.0.0", "--port", "8000"]}.
\end{enumerate}
\code{.dockerignore} excludes \code{.env}, \code{index.pkl},
\code{__pycache__/}, \code{*.pyc}, \code{.venv/}, \code{venv/}, \code{.git/},
\code{.gitignore}, and all \code{*.md} files --- so local secrets and index
data are not baked into the image, and configuration is supplied at runtime.
Build and run:
\begin{lstlisting}[language=bash, style=bashstyle]
docker build -t rag-project .
docker run -d --name rag-app --restart unless-stopped -p 8000:8000 \
  --env-file .env -v "$(pwd)/data":/app/data rag-project
\end{lstlisting}
For the volume mount to persist data, \code{.env} must set
\code{INDEX\_PATH=/app/data/index.pkl} and
\code{DOCS\_FOLDER=/app/data/sample\_docs}.

% =================================================================
\section{Deployment}
Two deployment paths exist, using identical code; only
\code{GENERATION\_BACKEND} differs.

\subsection{AWS EC2 with a Python Venv (Gemini)}
Documented in \code{DEPLOY\_AWS.md}. A single EC2 instance (Ubuntu 22.04,
\texttt{t3.small}) runs uvicorn under \texttt{systemd} (\code{rag.service})
for automatic restart on boot or crash. Gemini supplies generation. Covers
launching the instance, installing dependencies, adding the API key, service
registration, verification, and optional hardening (Elastic IP, nginx on port
80, HTTPS via certbot, and S3 for the index).

\subsection{Docker + Amazon Bedrock (Claude)}
Documented in \code{DEPLOY\_DOCKER\_BEDROCK.md}. The same EC2 host runs the
app in a Docker container; answers come from Amazon Bedrock using an IAM
instance role so that \emph{no API key} is managed. Covers the one-time
Anthropic use-case submission (if required), IAM role creation, Docker
installation, Bedrock-specific environment configuration, the exact run
flags, and update procedures.

% =================================================================
\section{Security and Privacy}
\begin{itemize}
  \item \code{.env} and \code{index.pkl} are gitignored and excluded from
        Docker images.
  \item The local \code{GEMINI\_API\_KEY} value does not match the usual
        Google AI Studio (\code{AIza...}) format and may be invalid;
        a fresh key from Google AI Studio is recommended if authentication
        errors occur.
  \item Generation errors are currently returned to the browser as the
        answer body (\code{ERROR generating answer: ...}), which can leak
        internal details; logging server-side and returning a generic message
        is advised before public deployment.
  \item One of the shipped sample documents
        (\code{sample\_docs/New Text Document.txt}) contains private
        personal-note content and should be removed before sharing or
        deploying this repository.
\end{itemize}

% =================================================================
\section{Known Limitations}
\begin{itemize}
  \item No OCR --- scanned or image-only PDFs yield no text.
  \item Fixed-size word chunking only; no structure-aware or semantic
        chunking.
  \item Retrieval is top-3 cosine similarity only --- no hybrid (keyword +
        vector) search, no reranker, no metadata filtering.
  \item The vector store is an in-memory numpy array persisted with pickle:
        no incremental merge on load, no ANN index (HNSW/IVF), no horizontal
        scaling, no transactional safety.
  \item No automated test suite or evaluation set; quality is assessed
        manually via retrieval scores and answer inspection.
  \item Single-process, single-user: no authentication, rate limiting,
        multi-user separation, or locking on index writes.
\end{itemize}

% =================================================================
\section{Future Work}
\begin{enumerate}
  \item Stress-test the prompt: verify the model says \emph{I don't know}
        rather than hallucinating when the context lacks an answer.
  \item Tune chunk size (e.g., 100 and 800) and observe changes in retrieval
        scores and answer quality.
  \item Add more documents and observe when 3-chunk retrieval stops finding
        the right passage; that signals the need for hybrid search or a
        reranker.
  \item Replace the numpy store with Chroma once \code{index.pkl} feels
        limiting; the interface is small enough to swap cleanly.
  \item Build an evaluation set of 15--20 questions with known answers and
        gate every change on pass/fail results.
\end{itemize}

% =================================================================
\section{Conclusion}
The project demonstrates a complete, transparent RAG pipeline in roughly a
thousand lines of dependency-light Python. Every stage---chunking,
embedding, retrieval, and generation---is inspectable and interchangeable,
which makes it a strong starting point for understanding RAG fundamentals and
for controlled experimentation before adopting heavier infrastructure.
Both a browser UI and a CLI surface the same pipeline, and generation can be
switched between Google Gemini and Amazon Bedrock at runtime without code
changes.

\end{document}